\uuid{68Pw}
\titre{Valeur principale de $1/x$}

\niveau{L3}
\module{Analyse}
\chapitre{Distributions}
\sousChapitre{Valeur principale et dérivation distributionnelle}
\theme{Dérivée de $\ln|x|$}
\auteur{Quentin Liard}
\datecreate{ 2026-09-03}
\organisation{CentraleSupélec}
\difficulte{4}

\contenu{

\texte{
On considère la distribution valeur principale, notée $\operatorname{vp}(1/x)$, définie par
$$
\left\langle \operatorname{vp}\left(\frac1x\right),\varphi\right\rangle
=
\lim_{\varepsilon\to0^+}
\int_{|x|>\varepsilon}\frac{\varphi(x)}{x}\,dx,
\qquad \forall\varphi\in\mathcal{D}(\mathbb{R}).
$$
}

\begin{enumerate}

\item \question{Montrer que la dérivée de $x\mapsto\ln|x|$ au sens des distributions est $\operatorname{vp}(1/x)$.}

\reponse{
La fonction $x\mapsto\ln|x|$ appartient à $L^1_{\mathrm{loc}}(\mathbb{R})$ et définit donc une distribution régulière. Soit $\varphi\in\mathcal{D}(\mathbb{R})$. On a
$$
\left\langle (\ln|x|)',\varphi\right\rangle
=-\int_{\mathbb{R}}\ln|x|\,\varphi'(x)\,dx.
$$
Pour $\varepsilon>0$, posons
$$
I_\varepsilon
=
\int_{-\infty}^{-\varepsilon}\ln(-x)\varphi'(x)\,dx
+
\int_{\varepsilon}^{+\infty}\ln(x)\varphi'(x)\,dx.
$$
Une intégration par parties sur chacun des deux intervalles donne
$$
I_\varepsilon
=
\ln(\varepsilon)\bigl(\varphi(-\varepsilon)-\varphi(\varepsilon)\bigr)
-
\int_{|x|>\varepsilon}\frac{\varphi(x)}{x}\,dx.
$$
Écrivons $\varphi(x)=\varphi(0)+x\psi(x)$, avec $\psi\in C^\infty(\mathbb{R})$. Alors
$$
\varphi(-\varepsilon)-\varphi(\varepsilon)
=-\varepsilon\bigl(\psi(-\varepsilon)+\psi(\varepsilon)\bigr),
$$
et donc
$$
\ln(\varepsilon)\bigl(\varphi(-\varepsilon)-\varphi(\varepsilon)\bigr)
\longrightarrow 0.
$$
Par conséquent,
$$
-I_\varepsilon
\longrightarrow
\left\langle \operatorname{vp}\left(\frac1x\right),\varphi\right\rangle.
$$
Or
$$
\left\langle (\ln|x|)',\varphi\right\rangle
=-\lim_{\varepsilon\to0^+}I_\varepsilon.
$$
Ainsi,
$$
\boxed{(\ln|x|)'=\operatorname{vp}\left(\frac1x\right)}
$$
dans $\mathcal{D}'(\mathbb{R})$.
}

\item \question{Montrer que
$$
x\,\operatorname{vp}\left(\frac1x\right)=1
$$
au sens des distributions.}

\reponse{
Soit $\varphi\in\mathcal{D}(\mathbb{R})$. Par définition du produit d'une distribution par une fonction lisse,
$$
\left\langle x\,\operatorname{vp}\left(\frac1x\right),\varphi\right\rangle
=
\left\langle \operatorname{vp}\left(\frac1x\right),x\varphi\right\rangle.
$$
Ainsi,
$$
\left\langle x\,\operatorname{vp}\left(\frac1x\right),\varphi\right\rangle
=
\lim_{\varepsilon\to0^+}
\int_{|x|>\varepsilon}\frac{x\varphi(x)}{x}\,dx
=
\lim_{\varepsilon\to0^+}
\int_{|x|>\varepsilon}\varphi(x)\,dx.
$$
Comme $\varphi$ est intégrable,
$$
\lim_{\varepsilon\to0^+}
\int_{|x|>\varepsilon}\varphi(x)\,dx
=
\int_{\mathbb{R}}\varphi(x)\,dx
=
\langle 1,\varphi\rangle.
$$
On en déduit
$$
\boxed{x\,\operatorname{vp}\left(\frac1x\right)=1}.
$$
}

\end{enumerate}

}
